\documentclass[twocolumn,pra,superscriptaddress,nofootinbib]{revtex4-2}

\usepackage{amsmath}
\usepackage{amssymb}
\usepackage{amsthm}
\usepackage{mathtools}
\usepackage{booktabs}
\usepackage{hyperref}
\usepackage{array}
\usepackage{enumitem}
\usepackage{xcolor}

%% Theorem environments
\newtheorem{theorem}{Theorem}
\newtheorem{proposition}{Proposition}
\newtheorem{lemma}{Lemma}
\newtheorem{corollary}{Corollary}
\newtheorem{conjecture}{Conjecture}
\newtheorem{definition}{Definition}
\newtheorem{observation}{Observation}
\newtheorem{remark}{Remark}

%% Commands
\newcommand{\R}{\mathbb{R}}
\newcommand{\C}{\mathbb{C}}
\newcommand{\Z}{\mathbb{Z}}
\newcommand{\Q}{\mathbb{Q}}
\newcommand{\KS}{\mathrm{KS}}
\newcommand{\CK}{\mathrm{CK\text{-}31}}

\begin{document}

\title{The 24--31 gap: Computational methods and structural barriers for minimal Kochen--Specker sets in dimension three}

\author{Michael Kernaghan}
\affiliation{Pacific Quantum Systems, Vancouver, Canada}

\date{March 2026}

\begin{abstract}
The minimum number of rays required for a Kochen--Specker set in dimension~3
is bounded between 24 and~31, a gap that has persisted for over a decade.
We report a systematic search for KS sets with fewer than~31 vectors using
seven complementary strategies: SAT-verified minimization within six algebraic
ray pools, cross-pool mixing, expanded coordinate alphabets (30~tested, including
Pisot-number fields), numerical optimization, deletion-criticality analysis,
triad density thresholds, and a graph perturbation search over 500,000 trials.
No sub-31 KS set is found by any strategy.
An OCUS procedure certifies that~31 is the exact minimum within the 49-ray
integer pool $\{0,\pm1,\pm2\}^3$.
The graph perturbation search finds 61,702 abstract KS-uncolorable hypergraphs
with fewer than~31 vertices; numerical realizability testing finds no
$\R^3$ embedding for any of them (best residual $3.8 \times 10^{-2}$).
A complementary vertex-merging analysis produces 394 combinatorially valid
30-vertex KS graphs, all proved unrealizable within the integer pool by
finite SAT encoding.
This identifies a \emph{realizability barrier}: abstract sub-31 KS candidates
exist in abundance, but three-dimensional geometry prevents their realization.
We establish an empirical \emph{modulus-2 boundary}---among integer-like
alphabets $\{0,\pm1,\pm x\}$, only those where $|x|^2 = 2$ (or the
degenerate case $x = 2$) produce KS-uncolorable pools---and
analyze structural constraints (constraint budget, plane matching property)
relevant to proof attempts.
The paper surveys prior work on KS minimality, taxonomizes the available
methods of attack, and maps the landscape of approaches for closing the
24--31 gap.
\end{abstract}

\maketitle

%%----------------------------------------------------------------------
\section{Introduction}
\label{sec:intro}
%%----------------------------------------------------------------------

The Kochen--Specker (KS) theorem \cite{KochenSpecker1967} (see Ref.~\cite{Budroni2022} for a comprehensive review) establishes that quantum
measurements in Hilbert spaces of dimension $d \geq 3$ cannot be described by
noncontextual hidden-variable theories. Formally, the theorem is witnessed by a
finite set of rays in $\R^d$ (equivalently, unit vectors up to sign) together with
a collection of mutually orthogonal triples---the \emph{bases} or \emph{contexts}---such
that no assignment of values in $\{0,1\}$ to the rays can simultaneously satisfy two
conditions: (i) exactly one ray in each basis receives the value~1, and (ii) if two
rays from distinct bases are collinear, they receive the same value. Any such
configuration is called a \emph{KS set}. Its existence rules out hidden-variable
completions that assign pre-measurement values independently of context.

Determining the minimum number of rays required to form a KS set in dimension~3
is one of the central open problems in quantum foundations. Despite more than
five decades of effort, a definitive answer remains out of reach. The current
situation is captured by the inequality $24 \leq n_{\min}^{(3)} \leq 31$, and
closing this gap is the motivating problem of the present paper.

This paper is a companion to Ref.~\cite{Kernaghan2026algebraic}, which
establishes the algebraic framework---six discrete ``algebraic islands'' of
KS-supporting number rings, two cancellation mechanisms, and the pool
construction methodology---on which the present search is built. We assume
that framework here and do not re-derive it; the reader is referred to
Ref.~\cite{Kernaghan2026algebraic} for definitions, proofs, and the full
classification. The present paper's contribution is the systematic search
for sub-31 KS sets using that framework, the identification of structural
barriers, and the mapping of approaches for closing the gap.

\subsection*{History of the upper bound}

The first explicit KS proof in dimension~3, due to Kochen and Specker
\cite{KochenSpecker1967}, used 117 rays. Over the following decades, that number
fell steadily. Peres~\cite{Peres1991} provided a 33-vector proof drawing rays
from the coordinate alphabet $\{0, \pm 1, \pm\sqrt{2}\}$, and later reported the
Conway--Kochen 31-vector set (CK-31) in his textbook \cite{Peres1993,ConwayKochen}.
The CK-31 construction works entirely within the integer alphabet
$\{0,\pm 1, \pm 2\}$: exact orthogonalities arise from the
cancellation identity $1^2 + 1^2 = (\sqrt{2})^2$, which in integer coordinates
reads $1 + 1 = 2$ after squaring. This arithmetic mechanism produces a remarkably
compact combinatorial structure---31 rays arranged in 17 orthogonal
triples (bases)---that has resisted all attempts at reduction for over three decades.
The CK-31 set therefore stands as the smallest \emph{known} KS set in dimension~3.

Cabello's recent construction \cite{Cabello2025simplest}, using complex coordinates
drawn from the Eisenstein integers $\Z[\omega]$ with $\omega = e^{2\pi i/3}$,
yields 33 rays and provides an independent upper bound at~33, distinguished by
minimizing the number of bases (14) rather than the number of rays. While it does
not improve the upper bound on $n_{\min}^{(3)}$, it demonstrates that algebraic
fields beyond $\Z$ can support small KS sets and illuminates the role of complex
structure in KS geometry.

\subsection*{History of the lower bound}

Progress on the lower bound has been slower but accelerating. Arends, Ouaknine,
and Wampler~\cite{ArendsOuaknineWampler2011} established a lower bound of~18 via
systematic combinatorial search. Uijlen and Westerbaan~\cite{UijlenWesterbaan2016}
raised this to~22 using a SAT-based approach with certified uncolorability.
The bound advanced to~24 independently and concurrently: Li, Bright, and
Ganesh~\cite{LiBrightGanesh2024} used SAT solvers augmented with algebraic
realizability checking, while Kirchweger, Peitl, and
Szeider~\cite{KirchwegarPeitlSzeider2023} employed co-certificate learning with
SAT modulo symmetries. Both works confirm that no KS set on fewer than~24 rays
exists in $\R^3$ or $\C^3$. The gap between~24 and~31 is the current frontier.

\subsection*{Rigidity and the conjecture that 31 is optimal}

Trandafir and Cabello~\cite{TrandafirCabello2025} recently proved that CK-31 is
\emph{rigid}: it is unique up to unitary equivalence in $\C^3$, and every
sub-configuration of fewer than~31 rays drawn from its basis structure is
KS-colorable. They further conjectured, subject to a stated structural assumption,
that $n_{\min}^{(3)} = 31$. This conjecture represents the sharpest current
theoretical statement about the problem and provides the primary motivation for
the present investigation: do the methods available today support or refute it?

\subsection*{Our contribution}

We report a systematic, multi-strategy computational search for KS sets with
fewer than~31 vectors. Seven complementary methods are deployed, spanning
algebraic enumeration, SAT-based optimization, numerical gradient descent,
graph-theoretic perturbation, and structural analysis. No sub-31 KS set is
found by any strategy.

Our main findings are as follows. First, within each of six algebraic ray pools
identified in our companion paper \cite{Kernaghan2026algebraic}, MUS (minimum
unsatisfiable subset) extraction consistently finds sharp minima: 31 rays for the
integer pool, 33 for Peres, Eisenstein, and $\Z[\sqrt{-2}]$, 43 for Heegner-7,
and 52 for the golden ratio. For the integer pool, an OCUS (optimal constrained
unsatisfiable subset) procedure exhaustively certifies that no KS-uncolorable
subset of $\leq 30$ rays exists among all 49 pool rays; for the remaining pools,
the minima are best-known values supported by thousands of independent MUS
extractions but not exhaustively certified.

Second, we identify an empirical \emph{modulus-2 boundary}: among
integer-like alphabets of the form $\{0,\pm 1,\pm x\}$, only those where
$|x|^2 = 2$ produce KS-uncolorable pools. Alphabets with $|x|^2 \geq 3$
(such as $\{0,\pm 1,\pm 3\}$) are colorable despite generating orthogonal
triples. In a companion paper~\cite{Kernaghan2026algebraic}, we established
that KS-uncolorability in dimension~3 requires one of two cancellation
mechanisms: modulus-2 cancellation ($|x|^2 = 2$, of which the integer
identity $1+1=2$ is the degenerate case) or phase cancellation (a vanishing
sum of roots of unity, as in the Eisenstein field). The modulus-2 boundary
reported here is the integer-alphabet manifestation of this broader pattern.

Third, a graph perturbation search over 500,000 random trials finds 61,702
abstract KS-uncolorable hypergraphs with fewer than~31 vertices; numerical
optimization finds no $\R^3$ realization for any of them (best residual
$3.8 \times 10^{-2}$, six orders of magnitude above the acceptance threshold).
A complementary vertex-merging analysis produces 394 distinct 30-vertex
KS-uncolorable graphs derived from CK-31; all 394 are proved unrealizable
within the integer pool by finite SAT encoding ($< 1$\,s each).
The convergence of these two independent lines of evidence---one numerical,
one exact within a finite pool---identifies a \emph{realizability barrier}:
abstract sub-31 KS candidates exist in abundance, but embedding them as
rays in $\R^3$ appears to be the core obstruction.

Fourth, we analyze the structural reasons why naive counting and coloring
arguments are insufficient to close the gap, and identify the \emph{plane
matching property}---the requirement that every ray participate in sufficiently
many orthogonal triples drawn from a fixed geodesically constrained
configuration---as a candidate geometric constraint for future proof attempts.

\subsection*{Organization}

Section~\ref{sec:prior} surveys prior work on KS minimality in detail.
Section~\ref{sec:methods} introduces our computational framework, including
the SAT encoding, realizability oracle, and software stack.
Sections~\ref{sec:exact}--\ref{sec:density} present six search strategies:
exact minimization within algebraic pools, cross-pool mixing, expanded alphabets,
numerical optimization, criticality analysis, and triad density thresholds.
Section~\ref{sec:perturbation} presents the graph perturbation search.
Section~\ref{sec:merging} analyzes vertex merging and the realizability barrier.
Section~\ref{sec:proof} develops structural insights toward a proof of optimality.
Sections~\ref{sec:discussion} and~\ref{sec:conclusion} discuss the landscape of
approaches and state the open problems that remain.

The paper is written as an exploration of what current methods can and cannot
achieve, addressed to researchers in quantum foundations, combinatorics, and
automated reasoning. The central question---whether $n_{\min}^{(3)} = 31$---remains
open, but the evidence assembled here points consistently in one direction.

%%----------------------------------------------------------------------
\section{Prior work on KS minimality}
\label{sec:prior}
%%----------------------------------------------------------------------

We organize the literature around four themes: the succession of upper bounds,
the progression of lower bounds, recent rigidity results, and the algebraic
framework that underlies our search.

%%- - - - - - - - - - - - - - - - - - - - - - - - - - - - - - - - - - -
\subsection{Upper bounds}
\label{sec:prior:upper}
%%- - - - - - - - - - - - - - - - - - - - - - - - - - - - - - - - - - -

The first explicit proof of the Kochen--Specker theorem in dimension~3 was given
by Kochen and Specker themselves in 1967~\cite{KochenSpecker1967}.  Their
construction used 117 rays arranged in a set of interlocking orthogonal triples
that admits no consistent $\{0,1\}$-coloring.  While the proof was conceptually
foundational, the sheer size of the configuration left open the question of how
compact a KS proof could be made.

The first significant reduction came from Peres~\cite{Peres1991}, who in 1991
exhibited two 33-vector proofs drawing coordinates from the alphabet
$\{0, \pm 1, \pm\sqrt{2}\}$.  The non-integer entry $\sqrt{2}$ was essential:
Peres exploited the identity $1^2 + 1^2 = (\sqrt{2})^2$ to produce exact
orthogonalities that cannot be achieved with integer coordinates alone.
In the same work, Peres noted that John~H.~Conway and Simon~Kochen had
communicated, around 1990, a 33-vector construction using only integer
coordinates (subsequently referred to as CK-33).  Both the Peres 33-vector
set and CK-33 contain~33 vectors; they differ in their alphabets and in
their orthogonality graphs.

Conway and Kochen subsequently improved their construction to 31 vectors, using
coordinates drawn exclusively from the integer alphabet $\{0, \pm 1, \pm 2\}$.
This CK-31 construction was reported by Peres in his 1993
textbook~\cite{Peres1993,ConwayKochen}.  The cancellation underlying CK-31 is
the identity $1 + 1 = 2$ (equivalently, $1^2 + 1^2 + (-1)(2) = 0$ in a suitable
triple product), which produces~17 mutually orthogonal triples from just~31 rays.
CK-31 has remained the record holder for over three decades; no subsequent work
has found a KS set in dimension~3 with fewer than~31 rays.

A notable observation about all upper bounds in dimension~3 is their reliance on
algebraic coordinate structure~\cite{Pavicic2004,Pavicic2019}.  Every known KS set in $\R^3$ with $\leq 33$
rays has coordinates drawn from a discrete, algebraically defined alphabet: the
integers $\{0,\pm 1,\pm 2\}$, the ring $\Z[\sqrt{2}]$, or a cyclotomic extension.
No ``generic'' or ``numerically generic'' KS set with an irregular coordinate
distribution has ever been found.  This empirical rigidity of the upper bound
motivates the algebraic approach taken in the present paper.

Cabello~\cite{Cabello2025simplest} recently introduced a 33-vector KS set
constructed from vectors over the Eisenstein integers $\Z[\omega]$ with
$\omega = e^{2\pi i/3}$, using a Weyl--Heisenberg symmetry to organize the
construction.  This set uses only~14 bases (orthogonal triples), fewer than any
previously known KS set in dimension~3, and was proposed as the ``simplest'' KS
set by the criterion of minimizing the number of measurement contexts.  It does
not improve the upper bound on the number of rays ($33 > 31$), but it establishes
that fields beyond the integers can produce small KS sets and that the number of
bases is a distinct minimization target.

For completeness, we mention record constructions in dimensions $d > 3$, which
demonstrate that the minimum number of rays required for a KS proof decreases
sharply with dimension.  Kernaghan~\cite{Kernaghan1994} found a 20-vector KS set
in dimension~4, and Kernaghan and Peres~\cite{KernaghanPeres1995} gave a 36-vector
construction in dimension~8.  Cabello, Estebaranz, and
Garc\'ia-Alcaine~\cite{CabelloEstebaranzGarcia1996} subsequently reduced the
dimension-4 record to 18 vectors.  These results are not directly comparable to
the dimension-3 problem: the combinatorial and algebraic constraints differ
fundamentally between dimensions, and a KS set for $d = 4$ does not embed into
$\R^3$.

%%- - - - - - - - - - - - - - - - - - - - - - - - - - - - - - - - - - -
\subsection{Lower bounds}
\label{sec:prior:lower}
%%- - - - - - - - - - - - - - - - - - - - - - - - - - - - - - - - - - -

Establishing that any KS proof in dimension~3 requires at least $n$ rays is
significantly harder than exhibiting a KS set.  Early lower bounds were modest:
simple counting and coloring arguments show that a KS set in dimension~3 needs
at least~5 rays, and incremental combinatorial improvements over several decades
brought this to~15.

Arends, Ouaknine, and Wampler~\cite{ArendsOuaknineWampler2011} raised the lower
bound to~18 in 2011 using a systematic approach combining exhaustive enumeration
of candidate orthogonality hypergraphs with SAT-based coloring checks.  Their
method established that every 3-uniform hypergraph on fewer than~18 vertices whose
hyperedges form a system of mutually orthogonal triples admits a valid $\{0,1\}$
KS-coloring in $\R^3$.

Uijlen and Westerbaan~\cite{UijlenWesterbaan2016} improved the bound to~22 in
2016 using stronger combinatorial constraints on the orthogonality graph.  Their
approach incorporated more refined structure on the allowed hypergraph types,
ruling out additional families of small candidate configurations.

Two independent groups then established the current best lower bound of~24 vectors.

Li, Bright, and Ganesh~\cite{LiBrightGanesh2024}, in work presented at IJCAI 2024,
developed a pipeline that works in three stages: (1)~enumerate abstract 3-uniform
hypergraphs on $n$ vertices that are KS-uncolorable (as verified by a SAT solver);
(2)~filter these by checking, again with a SAT solver, that the abstract
uncolorability cannot be destroyed by any assignment of triples; and (3)~test
whether each surviving abstract hypergraph can be \emph{geometrically realized}
as actual rays in $\R^3$ by submitting the embeddability query to Z3's nonlinear
arithmetic solver.  The geometric realizability step is the computational
bottleneck: Z3 returns ``unknown'' on embeddability queries for hypergraphs of
order greater than approximately~20, because the resulting systems of polynomial
equations exceed the solver's decision capabilities.  Despite this bottleneck, the
pipeline is rigorous: any abstract hypergraph for which Z3 certifies
non-realizability is definitively ruled out.  Li, Bright, and Ganesh concluded
that no KS set on fewer than~24 rays exists in $\R^3$.

Kirchweger, Peitl, and Szeider~\cite{KirchwegarPeitlSzeider2023}, in independent
work presented at AAAI 2023, obtained the same lower bound of~24 using a different
technique: co-certificate learning within a SAT modulo symmetries framework.  Their
approach generates symmetry-breaking clauses automatically to prevent the
enumeration from revisiting isomorphic configurations, and uses co-certificates
(proofs of unsatisfiability) to certify that each explored configuration is
genuinely unrealizable.  The method does not rely on external solvers for
realizability and provides a self-contained proof.

Both groups confirm that no KS set exists on fewer than~24 rays in $\R^3$ or
$\C^3$.  Neither group claims that~24 is tight---the actual minimum $n_{\min}^{(3)}$
may lie anywhere in the interval $[24, 31]$.  Closing this gap remains the
central open problem in KS minimality, and the primary challenge is the
realizability bottleneck: abstract KS-uncolorable hypergraphs with far fewer
than~31 vertices exist in abundance, but verifying whether any can be embedded
as real or complex rays requires algebraic tools that currently do not scale to
the problem.

%%- - - - - - - - - - - - - - - - - - - - - - - - - - - - - - - - - - -
\subsection{Rigidity and uniqueness}
\label{sec:prior:rigid}
%%- - - - - - - - - - - - - - - - - - - - - - - - - - - - - - - - - - -

A KS set is \emph{rigid} if every other set of rank-1 projectors in $\C^3$ whose
orthogonality graph is isomorphic to that of the original must be unitarily
equivalent to it.  Rigidity is a strong uniqueness statement: it implies that the
geometric realization is essentially unique, up to the symmetries of Hilbert space.

Trandafir and Cabello~\cite{TrandafirCabello2025} recently proved that CK-31 is
rigid in $\C^3$.  Specifically, they showed that any collection of $31$ rank-1
projectors satisfying the same orthogonality relations as CK-31---i.e.,
any geometric realization of the CK-31 orthogonality graph---must be unitarily
equivalent to the original Conway--Kochen construction.  This means CK-31 admits
no non-trivial deformations: it cannot be continuously or discretely modified
while preserving its combinatorial structure.  In the same work, they proved
the analogous rigidity result for CK-33 (the earlier Conway--Kochen
33-vector set, which is distinct from both CK-31 and the Peres 33-vector
construction).

The rigidity result has a direct implication for the minimality problem.
Because CK-31 is rigid, any KS set with fewer than~31 rays must have an
orthogonality graph that is \emph{not isomorphic} to that of CK-31.  It cannot
arise by a minor modification, perturbation, or deformation of the existing
31-ray construction: a sub-31 KS set, if it exists, would require an entirely
different combinatorial structure.

Trandafir and Cabello also formulated a conditional conjecture:
subject to a structural assumption about the relationship between abstract
KS-uncolorable hypergraphs and their geometric realizations in $\C^3$, they
conjectured that $n_{\min}^{(3)} = 31$.  This conjecture is the sharpest
theoretical statement currently available about KS minimality in dimension~3
and provides the principal motivation for the computational investigation
reported in the present paper.

%%- - - - - - - - - - - - - - - - - - - - - - - - - - - - - - - - - - -
\subsection{The algebraic framework}
\label{sec:prior:algebraic}
%%- - - - - - - - - - - - - - - - - - - - - - - - - - - - - - - - - - -

All known small KS sets in dimension~3 exploit specific algebraic structures in
their coordinate systems.  Exact orthogonality between two integer-coordinate
rays requires a vanishing inner product over $\Z$, which in turn requires the
existence of a \emph{cancellation identity}: an algebraic relation among
coordinate values that allows their pairwise products to sum to zero.  The
abundance or scarcity of such relations determines how many orthogonal triples
a given coordinate alphabet can support, and therefore how small a KS set
constructed from that alphabet can be.

In a companion paper~\cite{Kernaghan2026algebraic}, we carry out a systematic
survey of coordinate alphabets arising from quadratic, cyclotomic, and
golden-ratio number fields, identifying which fields support KS constructions
in dimension~3 and characterizing the algebraic mechanism in each case.

The survey identifies six \emph{discrete algebraic islands}: disjoint families
of coordinate alphabets, each internally consistent and each supporting a
distinct smallest KS size.  The islands correspond to (i)~the rational integers
$\Z$ with alphabet $\{0,\pm 1,\pm 2\}$ (minimum 31 rays), (ii)~the Eisenstein
integers $\Z[\omega]$ (minimum 33), (iii)~the ring $\Z[\sqrt{2}]$ associated
with Peres (minimum 33), (iv)~the Gaussian-like ring $\Z[\sqrt{-2}]$ (minimum
33), (v)~the Heegner-7 ring $\Z[\tfrac{1+\sqrt{-7}}{2}]$ (minimum 43), and
(vi)~the golden-ratio ring $\Z[\varphi]$ with $\varphi = \tfrac{1+\sqrt{5}}{2}$
(minimum 52).  These six families are the only algebraic sources of KS
constructions in dimension~3 found within the search space.

Two distinct cancellation mechanisms account for all observed KS-uncolorability.
The first is \emph{modulus-2 cancellation}: the generating algebraic integer $x$
satisfies $|x|^2 = 2$, allowing the inner-product identity $x\bar{x} = 2$
to produce exact orthogonal triples.  The pure integer case, $1 + 1 = 2$, is
the degenerate additive instance of this mechanism.  The second is
\emph{phase cancellation}: a vanishing sum of roots of unity, of which the
simplest example is $1 + \omega + \omega^2 = 0$ for $\omega = e^{2\pi i/3}$,
underlies the Eisenstein island.  All six algebraic islands fall under one of
these two mechanisms.  Coordinate alphabets in which neither mechanism is
available---including integer alphabets such as $\{0, \pm 1, \pm 3\}$ that lack
the entry $\pm 2$---produce only colorable ray pools in all tested cases.

The companion paper also introduces two new KS constructions not previously
reported in the literature: a 43-vector set based on the Heegner-7 ring and a
52-vector set based on the golden-ratio ring.  These new constructions confirm
the existence of the corresponding algebraic islands and demonstrate that KS
geometry is richer than a focus on integer coordinates alone would suggest.

The present paper builds directly on this algebraic framework.  Rather than
searching over arbitrary ray configurations, we restrict attention to the six
identified algebraic pools and ask whether any of them supports a KS set with
fewer than~31 vectors, using the algebraic structure to make the search tractable.

%%----------------------------------------------------------------------
\section{Framework and methods}
\label{sec:methods}
%%----------------------------------------------------------------------

\subsection{Algebraic pool construction}

A \emph{coordinate alphabet} $\mathcal{A}$ is a finite set of values from which ray
coordinates are drawn. The \emph{induced ray pool} $S(\mathcal{A})$ consists of all
projectively distinct nonzero vectors in $\mathcal{A}^3$. Rays are canonicalized: the
first nonzero coordinate is taken positive, integer coordinates are $\gcd$-reduced to
primitive form, and complex coordinates are placed in a standard lexicographic normal
form. The \emph{orthogonality graph} of a pool has rays as vertices and edges between
orthogonal pairs (dot product zero); \emph{triads} are 3-cliques in this graph---mutually
orthogonal triples that form a basis up to normalization. See
Ref.~\cite{Kernaghan2026algebraic} for the full algebraic treatment of pool construction,
canonicalization, and the ring-theoretic conditions on $\mathcal{A}$ that determine
whether $S(\mathcal{A})$ contains any triads.

The six algebraic pools studied in this series are summarized in
Table~\ref{tab:pools}. Each arises from a distinct number ring, and the pool sizes
and triad counts determine the difficulty of the associated SAT and optimization
problems.

\begin{table}[t]
\caption{The six algebraic ray pools. See Ref.~\cite{Kernaghan2026algebraic} for
construction details.}
\label{tab:pools}
\begin{tabular}{llrr}
\toprule
Pool & Ring & Rays & Triads \\
\midrule
Integer      & $\Z$              &  49 &  26 \\
Eisenstein   & $\Z[\omega]$      &  57 &  22 \\
Peres        & $\Z[\sqrt{2}]$    &  49 &  16 \\
$\Z[\sqrt{-2}]$ & $\Z[\sqrt{-2}]$ & 49 &  16 \\
Heegner-7    & $\Z[\alpha]$      & 145 &  42 \\
Golden ratio & $\Z[\varphi]$     & 205 & 166 \\
\bottomrule
\end{tabular}
\end{table}

\subsection{SAT encoding of KS-uncolorability}

KS-uncolorability is encoded as a propositional satisfiability (SAT) instance as
follows. Introduce a Boolean variable $b_v$ for each ray $v$ in the pool, with the
intended interpretation that $b_v = 1$ means $v$ is colored ``green'' (assigned
value~1). For each triad $\{a, b, c\}$, enforce an exactly-one constraint via the
at-least-one clause
\begin{equation}
  (b_a \vee b_b \vee b_c)
\end{equation}
together with three pairwise-exclusion clauses $(\neg b_\alpha \vee \neg b_\beta)$
for each pair $\alpha \neq \beta \in \{a,b,c\}$. For each orthogonal pair $(u,v)$
not sharing any triad in the pool, add the exclusion clause $(\neg b_u \vee \neg b_v)$.
The resulting propositional formula is unsatisfiable if and only if no valid KS
coloring of the ray set exists, i.e., the set is a KS set. This encoding is the
standard extension of the triad constraint to non-basis orthogonal pairs, and is
strictly stronger than requiring only the triad constraints. All SAT solving uses the
Glucose4 solver~\cite{AudemardSimon2018} via the PySAT interface.

\subsection{MUS and OCUS extraction}

A \emph{Minimal Unsatisfiable Subset} (MUS) of a propositional formula is a subset of
clauses that is itself unsatisfiable but becomes satisfiable upon removal of any single
clause. In the KS encoding, a MUS identifies a minimal subset of orthogonality and
triad constraints that already forces uncolorability. To extract MUSes at the
\emph{ray level}---identifying minimal KS-uncolorable ray subsets---we introduce a
selector literal $a_r$ for each ray $r$, and condition all clauses involving $r$ on
the assumption $a_r = \mathrm{true}$. The solver then produces an unsatisfiable core
over selector literals. This core is minimized by iterative deletion: we remove one
selector at a time and check whether the residual formula remains unsatisfiable; a
selector is discarded if its removal preserves unsatisfiability, retained otherwise.
After minimization, the active selectors identify a minimal KS-uncolorable ray subset.
Each independent MUS extraction (with randomized deletion order) may produce a
different minimal subset, enabling enumeration of distinct minimal KS sets within a
pool.

\emph{OCUS} (Optimal Constrained Unsatisfiable Subset) strengthens MUS extraction by
certifying the \emph{globally minimum} unsatisfiable subset size within the pool.
OCUS iteratively proves, for increasing $k = 1, 2, \ldots$, that no unsatisfiable
subset of $\leq k$ rays exists, until an unsatisfiable subset of size $k^*$ is found
and its minimality is certified. This provides an exhaustive lower bound on the minimum
KS subset size within a finite pool, at the cost of substantially greater computational
effort than single-MUS extraction.

\subsection{Cross-product completion}

For some algebraic fields, the raw alphabet $\mathcal{A}$ produces a ray pool with too
few triads for KS-uncolorability. \emph{Cross-product completion} expands the pool by
adding rays of the form $v_i \times v_j$ for each orthogonal pair $(v_i, v_j)$ already
in the pool, where $\times$ denotes the standard vector cross product in $\R^3$. The
cross product of two vectors whose coordinates lie in a number field $K$ also has
coordinates in $K$, so the expanded pool remains within the same algebraic family. The
closure step is iterated until no new rays are generated. This procedure is essential
for the golden-ratio pool: the raw alphabet $\{0, \pm 1, \pm \varphi\}$ produces too
few triads to be KS-uncolorable, while the cross-product-completed pool of 205 rays
with 166 triads contains a 52-ray KS island that is invisible without this closure step.
See Ref.~\cite{Kernaghan2026algebraic} for the full procedure and its ring-theoretic
justification.

\subsection{Realizability testing}

Combinatorial KS-uncolorability (established by SAT) is a necessary but not sufficient
condition for a ray set to exist as an actual physical configuration in $\R^3$: the
abstract orthogonality graph must additionally be \emph{realizable}, meaning there must
exist actual unit vectors in $\R^3$ whose pairwise dot products match the graph's edge
structure. We employ three complementary methods for realizability testing, each with
distinct strengths and limitations.

\textbf{Z3/SMT.} Encode orthogonality as polynomial equations over the reals: for each
required-orthogonal pair $(v_i, v_j)$, impose $\sum_k v_{i,k} v_{j,k} = 0$, subject to
the normalization constraints $\|v_i\| = 1$. The Z3 SMT solver~\cite{Z3} can decide
satisfiability of such systems exactly using quantifier elimination over nonlinear real
arithmetic (NRA). This approach is exact and correct, but the NRA decision procedure is
doubly exponential in the number of variables; in practice, Z3 returns ``unknown'' for
instances with more than approximately 20 free variables, making it inapplicable to most
of the candidate sets encountered in this work.

\textbf{Numerical optimization.} Minimize the orthogonality residual
$\sum_{(i,j) \in E} (\langle v_i, v_j \rangle)^2$ subject to $\|v_i\| = 1$ for all $i$,
using L-BFGS-B with analytically computed gradients and 10--20 random initializations.
A final residual below $10^{-10}$ is taken as evidence of realizability. This approach
is fast and scales to large instances, but provides only one-sided evidence: a small
residual indicates realizability (up to numerical precision), while a large residual
does not prove unrealizability, since the optimizer may have converged to a local
minimum.

\textbf{Finite pool SAT.} For realizability \emph{within a specific algebraic pool} of
$m$ candidate rays, introduce Boolean variables $x_{v,r}$ (vertex $v$ of the abstract
graph assigned to pool ray $r$), with three constraint families: (i)~each vertex is
assigned exactly one ray ($\sum_r x_{v,r} = 1$); (ii)~distinct vertices receive distinct
rays ($x_{v,r} + x_{w,r} \leq 1$ for $v \neq w$); and (iii)~if vertices $v$ and $w$
are required to be orthogonal, their assigned rays must be orthogonal in the pool.
This reduces realizability within the pool to pure Boolean satisfiability, solvable
in under 1~s for the instances that defeat Z3. The trade-off is scope: the method tests
realizability within a finite algebraic pool, not in $\R^3$ generally. A negative
answer (UNSAT) means the abstract graph has no realization within that pool; it does
not rule out realizations in other pools or over the reals.

%%----------------------------------------------------------------------
\section{Minimization within algebraic pools}
\label{sec:exact}
%%----------------------------------------------------------------------

For each of the six algebraic ray pools (Table~\ref{tab:pools}), we applied
three heuristic minimization methods in concert with the SAT encoding of
Section~\ref{sec:methods}: (a)~clause-level MUS extraction with randomized
deletion, (b)~greedy ray deletion with 2000 random orderings plus deterministic
orderings (ascending and descending norm, ascending and descending degree), and
(c)~swap-based local search in which one ray at a time is replaced by a pool
ray not currently in the set, with the swap accepted whenever KS-uncolorability
is preserved. All three methods use Glucose4 as the underlying SAT oracle.

\subsection{Per-pool minima}

Table~\ref{tab:minima} summarizes the best result found per pool. For the
integer pool, all 5000 independent MUS extractions returned a set of exactly 31
rays, and no greedy or swap trial ever produced a 30-ray KS-uncolorable subset.
For the Eisenstein, Peres, and $\Z[\sqrt{-2}]$ pools, all 1000 MUS extractions
per pool returned sets of exactly 33 rays. The Heegner-7 pool minimum is 43 and
the golden-ratio pool minimum is 52.

\begin{table}[t]
\centering
\caption{Minimum KS subset size per pool. ``MUS min'' is the minimum found
across all MUS extractions. ``Trials'' is the number of independent MUS
extractions performed. Greedy and swap searches confirmed these minima in
every case.}
\label{tab:minima}
\begin{tabular}{lrrrr}
\toprule
Pool & Rays & Triads & MUS min & Trials \\
\midrule
Integer & 49 & 26 & 31 & 5000 \\
Eisenstein & 57 & 22 & 33 & 1000 \\
Peres & 49 & 16 & 33 & 1000 \\
$\Z[\sqrt{-2}]$ & 49 & 16 & 33 & 1000 \\
Heegner-7 & 145 & 42 & 43 & 1000 \\
Golden ratio & 205 & 166 & 52 & 1000 \\
\bottomrule
\end{tabular}
\end{table}

Greedy deletion (2000 random orderings plus norm-sorted and degree-sorted
deterministic orderings) and swap-based local search were run after each MUS
extraction as a consistency check: in no case did either method locate a subset
smaller than the MUS minimum. The convergence of three independent methods to
the same minimum provides strong corroborating evidence that the reported
values are true pool minima, though exhaustive proof is provided only by the
OCUS procedure described below.

\subsection{Structure of the six minimal 31-sets from the integer pool}

The following structural analysis was first reported in our companion
paper~\cite{Kernaghan2026algebraic}; we summarize it here because the norm
stratification and swap-isolation properties are directly relevant to the
sub-31 search.

From 5000 independent MUS extractions on the integer pool, exactly six
combinatorially distinct minimal 31-ray KS sets were recovered. (Two sets are
considered distinct if they differ in at least one ray; we cannot rule out the
existence of additional minimal 31-sets not reached by our randomized MUS
extractions.) Together, the six sets use 37 of the 49 available integer pool
rays, leaving 12 rays unused by any minimal set.

The six sets exhibit sharp norm stratification with respect to $\|v\|^2 =
v_1^2 + v_2^2 + v_3^2$:
\begin{itemize}
\item $\|v\|^2 \leq 3$ (the $\{0,\pm1\}$ sub-alphabet skeleton, 13~rays):
  present in \emph{all} six sets. These 13 rays form an invariant core.
\item $\|v\|^2 = 5$ (12~rays): present in 5 of 6 sets.
\item $\|v\|^2 = 6$ (12~rays): present in 4 of 6 sets.
\item $\|v\|^2 = 9$ (highest-norm rays in the pool, 12~rays): present in
  \emph{none} of the six sets.
\end{itemize}

The exclusion of all 12 highest-norm rays from every minimal 31-set is a
striking structural feature: the $\{0,\pm1,\pm2\}$ alphabet places the highest
norms at $\|v\|^2 = 9$ (e.g., $(2,2,1)$ type), and these rays are never
needed once the lower-norm rays supply sufficient interlocking constraints. The
invariant 13-ray core (all rays with $\|v\|^2 \leq 3$, i.e., rays using only
$\{0,\pm1\}$ coordinates) contributes 7 of the 26 integer-pool triads and forms
a configuration that is itself colorable, but whose extension by a small number
of norm-5 and norm-6 rays creates the interlocking necessary for
uncolorability.

Pairwise comparison of the six minimal sets reveals that no two are connected by
a single-ray swap: the minimum Hamming distance (number of ray replacements
needed to transform one set into another) is 5~ray replacements. Despite this
topological separation, all pairs share between 25 and 26 rays out of 31,
yielding a Jaccard similarity of approximately 0.71. CK-31, the Conway--Kochen
construction, is one of the six.

\subsection{OCUS exhaustive certification}

Heuristic methods cannot certify the absence of sub-31 sets; for this we used
the OCUS procedure. Starting from the full 49-ray integer pool and its SAT
encoding, OCUS iteratively enumerates hitting sets of all MUSes found so far,
calling the SAT oracle to test each candidate hitting set for unsatisfiability.
If a hitting set of size $k$ is UNSAT, a $k$-ray KS set has been found; if
no hitting set of size $k$ is UNSAT, OCUS advances to size $k{+}1$ or certifies
that no $k$-ray KS set exists within the pool.

Applied to the integer pool, OCUS exhausted all candidates of size $\leq 30$:
the search completed in 272 iterations, requiring less than 0.2~seconds of
wall-clock time on a single core. To provide a complete certificate across all
sizes from the trivial lower bound upward, we also ran OCUS for all values of
$n$ from 21 to 30 individually; the full range was certified in a combined
0.5~seconds. The result is an exhaustive, proof-by-exhaustion certificate:

\begin{quote}
\textit{No KS-uncolorable subset of $\leq 30$ rays exists within the full
49-ray integer pool $\{0,\pm1,\pm2\}^3$.}
\end{quote}

This rules out any sub-31 KS set constructible from the $\{0,\pm1,\pm2\}$
integer pool. Combined with the consistency across all six pools reported in
Table~\ref{tab:minima}, we conclude: \textit{no sub-31 KS set is constructible
from any of the six known algebraic ray pools in isolation.}

%%----------------------------------------------------------------------
\section{Cross-pool mixing}
\label{sec:cross}
%%----------------------------------------------------------------------

The per-pool minima of Section~\ref{sec:exact} establish that no single
algebraic pool yields a sub-31 KS set. However, this does not preclude the
possibility that rays from distinct pools might combine to create a smaller KS
set through \emph{cross-pool orthogonalities}: orthogonal pairs where one ray
belongs to pool~$A$ and the other to pool~$B$, mediated by a shared algebraic
sub-structure. This section tests that possibility.

\subsection{Method}

We formed all 15 pairwise pool combinations and the full union of all six pools.
For each combination, we: (1)~took the union of the two (or more) ray sets;
(2)~deduplicated rays by canonical form (first nonzero coordinate positive,
$\gcd$-reduced); (3)~computed the full orthogonality graph of the combined ray
set, including all cross-pool orthogonal pairs; and (4)~ran 500~independent
greedy minimization trials, each starting from a random ordering of the combined
ray set and deleting rays greedily while preserving KS-uncolorability. For each
combined pool, we recorded the minimum KS-uncolorable subset size found and the
pool composition of the resulting minimal set.

\subsection{Results}

Table~\ref{tab:cross} summarizes results for selected combinations; all 15
pairwise results are consistent with the pattern described here. The 477-ray
full union result is also included.

Despite significant cross-pool orthogonalities---for example, the Integer
$+$~Golden combination shares 78 cross-pool triads in which one or more rays
from each pool participate together in an orthogonal triple---every minimized set
draws its rays from a \emph{single pool only}:
\begin{itemize}
\item Any combination that includes the Integer pool: minimum~31, achieved
  entirely by integer rays.
\item Any combination that excludes the Integer pool: minimum~33, achieved
  entirely by rays from one non-integer pool.
\item The full 477-ray union of all six pools: minimum~31, achieved entirely
  by integer rays.
\end{itemize}

In each case, the greedy minimization procedure converged to a single-pool
solution: cross-pool orthogonalities did not provide a pathway to a smaller
mixed-pool KS set. The Integer pool's minimum of~31 dominates all combinations
in which it appears because its 31-ray sets already saturate the minimum
achievable by any of the heuristic methods employed.

\begin{table}[t]
\centering
\caption{Cross-pool mixing results for selected combinations. ``Min'' is the
best result in 500 greedy trials. ``Source'' indicates the pool from which all
rays in the minimal set are drawn. ``Cross-triads'' counts orthogonal triples
involving rays from two or more pools.}
\label{tab:cross}
\begin{tabular}{lrrl}
\toprule
Combination & Min & Cross-triads & Source \\
\midrule
Integer $+$ Eisenstein & 31 & 14 & Integer \\
Integer $+$ Golden & 31 & 78 & Integer \\
Integer $+$ Heegner-7 & 31 & 31 & Integer \\
Eisenstein $+$ Peres & 33 & 8 & Eisenstein \\
Peres $+$ $\Z[\sqrt{-2}]$ & 33 & 12 & Peres \\
Eisenstein $+$ Golden & 33 & 0 & Eisenstein \\
All six pools (477 rays) & 31 & --- & Integer \\
\bottomrule
\end{tabular}
\end{table}

\subsection{Caveat: heuristic limitations}

Greedy deletion is biased toward large basins of attraction in the space of
minimal KS sets. A mixed-pool minimal set with a small basin of attraction could
in principle exist but be systematically missed by greedy methods. Exhaustive
cross-pool search---testing all subsets of a combined pool for
KS-uncolorability---is computationally infeasible for the larger combinations
(e.g., the 477-ray union has $2^{477}$ subsets). The OCUS procedure of
Section~\ref{sec:exact} is in principle applicable to cross-pool combinations,
but the integer pool OCUS already runs on a 49-ray universe; extending it to the
477-ray union would require a significantly larger search.

We conclude: \textit{greedy cross-pool mixing finds no hybrid KS set smaller
than the corresponding single-pool minimum, ruling out hybrid algebraic
constructions from known fields subject to the heuristic nature of the search.}

%%----------------------------------------------------------------------
\section{Expanded alphabets}
\label{sec:alphabets}
%%----------------------------------------------------------------------

The six algebraic pools of Table~\ref{tab:pools} are each generated by a
specific coordinate alphabet. A natural question is whether a \emph{richer}
alphabet---more coordinate values, or a larger generator bound in the same
field---produces a denser orthogonality structure in which a smaller KS set
can be found. This section reports a systematic survey of 30 expanded alphabets.

\subsection{Method}

For each expanded alphabet, we: (1)~enumerated all projectively distinct rays in
$\R^3$ (or $\C^3$) whose coordinates lie in the extended alphabet; (2)~computed
the orthogonality graph and triad list; (3)~verified KS-uncolorability of the
full pool (some alphabets are not KS-uncolorable at all); and (4)~ran
300--500 greedy minimization trials to find the smallest KS-uncolorable subset.
Alphabets were drawn from four families: (i)~extended integer alphabets
$\{0,\pm1,\ldots,\pm k\}$; (ii)~extended versions of each non-integer pool
(larger generator bound or additional generator); (iii)~mixed-field alphabets
combining generators from two distinct rings; and (iv)~Pisot-number fields
(plastic ratio, tribonacci constant, silver ratio) not included among the
original six pools.

\subsection{Results}

Table~\ref{tab:alphabets} presents results for representative alphabets;
the full 30-alphabet survey is consistent with the patterns described here.

\begin{table}[t]
\centering
\caption{Selected expanded alphabets. ``KS'' indicates that the pool is
KS-uncolorable (--- means colorable); ``Min'' is the best greedy result in
300--500 trials (--- means not applicable).}
\label{tab:alphabets}
\begin{tabular}{lrcl}
\toprule
Alphabet & Rays & Min & Note \\
\midrule
$\{0,\pm1,\pm2\}$ & 49 & 31 & Baseline (CK-31) \\
$\{0,\pm1,\pm2,\pm3\}$ & 145 & 31 & Absorbs to CK-31 \\
$\{0,\pm1,\pm2,\pm3,\pm4\}$ & 289 & 36 & Heuristic limit \\
$\{0,\pm1,\pm3\}$ & 49 & --- & Not KS \\
$\{0,\pm1,\pm4\}$ & 49 & --- & Not KS \\
Eisenstein $c{=}3$, $n{\leq}6$ & 345 & 31 & Absorbs to CK-31 \\
Peres $+ \pm2$ & 109 & 31 & Absorbs to CK-31 \\
Heegner-7 $+ \pm2$ & 253 & 31 & Absorbs to CK-31 \\
Golden $+ \pm2$ & 637 & 33 & No absorption \\
$\{0,\pm1,\pm2,\pm i\}$ & 127 & 31 & Mixed-field \\
\bottomrule
\end{tabular}
\end{table}

\subsection{The modulus-2 boundary}

The most striking pattern in Table~\ref{tab:alphabets} is a sharp boundary
defined by the presence or absence of the coordinate value~$\pm2$.

\textbf{Alphabets containing $\pm2$ converge to minimum~31.} Every alphabet
that contains the $\{0,\pm1,\pm2\}$ sub-alphabet as a subset converges, under
greedy minimization, to a 31-ray KS set whose rays are drawn entirely from the
integer pool:
\begin{itemize}
\item $\{0,\pm1,\pm2,\pm3\}$ with 145 rays $\to$ minimum~31 (integer rays).
\item $\{0,\pm1,\pm2,\pm i\}$ with 127 rays (mixed real-imaginary field)
  $\to$ minimum~31 (integer rays).
\item Eisenstein $c=3$, $n \leq 6$, with 345 rays $\to$ minimum~31 (integer
  rays).
\item Peres $+ \pm2$ (109 rays) and Heegner-7~$+ \pm2$ (253 rays) $\to$
  minimum~31 (integer rays).
\end{itemize}
In each case the addition of $\pm2$ to a non-integer pool causes greedy
minimization to ``absorb'' the solution back to the CK-31 structure in the
integer sub-pool. The 637-ray Golden~$+\,\pm2$ pool is the sole exception:
greedy minimization finds minimum~33, and the golden-ratio algebraic structure
prevents absorption to the integer sub-pool.

\textbf{Alphabets lacking $\pm2$ are not KS-uncolorable.} The alphabets
$\{0,\pm1,\pm3\}$ and $\{0,\pm1,\pm4\}$---which possess 49~rays and 10~triads
each, comparable in size to the baseline integer pool---are \emph{not
KS-uncolorable at all}: the full pool is 2-colorable. Despite the availability
of orthogonal triples, the triads cannot be interlocked in a way that forces
contradiction. Analogously, $\{0,\pm1,\pm2,\pm3,\pm4\}$ with 289~rays reaches
only minimum~36 under greedy minimization---higher than the integer pool minimum
despite a much larger ray count---suggesting that the additional high-norm rays
do not contribute usable orthogonal structure.

The mechanism is algebraic: the cancellation identity $1+1=2$ enables
constructions such as $(1,1,1) \perp (1,1,-2)$, providing a ``two-against-one''
triad type that is absent from alphabets like $\{0,\pm1,\pm3\}$. Without this
identity, the only triads available are of ``one-against-one'' type or require
larger coordinate magnitudes, and these do not produce a sufficiently dense
interlocking structure for KS-uncolorability in dimension~3.

\subsection{Pisot-number fields}

We tested three higher-degree algebraic fields associated with Pisot numbers:
the plastic ratio ($x^3 = x+1$, degree~3), the tribonacci constant
($x^3 = x^2+x+1$, degree~3), and the silver ratio ($x^2 = 2x+1$, degree~2).
All three produce only colorable pools: despite generating orthogonal triples
from their respective algebraic identities, no KS-uncolorable configuration
was found. This extends the modulus-2 boundary observation beyond quadratic
fields to degree-3 algebraic numbers: the Pisot-number cancellation identities
are not sufficient for KS-uncolorability in the three-dimensional setting
surveyed here.

\subsection{Summary}

Across all 30 tested alphabets, no sub-31 KS set was found. The results
establish an empirical modulus-2 boundary for integer-like alphabets: among
alphabets of the form $\{0,\pm 1,\pm x\}$, only those with $|x|^2 = 2$
produce KS-uncolorable pools. (Non-integer fields such as the Eisenstein
integers use a different mechanism---phase cancellation---and are not
covered by this specific boundary; see Ref.~\cite{Kernaghan2026algebraic}
for the full two-mechanism classification.) We conclude: \textit{expanded
coordinate sets from all tested fields fail to produce sub-31 KS sets.}

%%----------------------------------------------------------------------
\section{Numerical optimization}
\label{sec:numerical}
%%----------------------------------------------------------------------

Strategies 1--3 operate exclusively within finite algebraic ray pools: they can only find KS sets whose coordinates lie in the tested number rings. A natural question is whether continuous optimization, unconstrained by algebraic structure, can locate KS sets in parts of configuration space not accessible to algebraic methods. We tested four non-algebraic methods as a \emph{negative control}.

\textbf{Simulated annealing from scratch.} We initialized random unit vectors in $\R^3$ (at sizes $n = 28, 29, 30$) and in $\C^3$ (at sizes $n = 28, 30$), then ran $2 \times 10^5$ annealing steps with an objective function that counts the number of exactly orthogonal pairs (inner product below a threshold $\epsilon = 10^{-8}$). Results: at most 5--6 exact orthogonal pairs were achieved in $\R^3$; zero exact orthogonal pairs were achieved in $\C^3$; and zero triads (mutually orthogonal triples) were found in either case across all tested sizes.

\textbf{CK-31 perturbation.} We removed $m = 1, 2, 3$ rays at random from CK-31, replaced them with random unit vectors drawn uniformly from $S^2$, and tested whether the resulting set is KS-uncolorable. In 1000 independent attempts for each $m$, the best result achieved was preservation of 16 of CK-31's 17 triads, but all 17 were never simultaneously preserved. KS-uncolorability was never restored.

\textbf{Random orthogonal completion.} Starting with 5 random seed rays, we iteratively extended the set by finding unit vectors orthogonal to existing pairs via the cross product $w = u \times v / \|u \times v\|$. Candidate sets of 60 rays achieved at most 1 triad. No KS-uncolorable set was found.

\textbf{Soft-tolerance annealing.} We allowed near-orthogonality (tolerance $\epsilon = 0.1$), then progressively tightened $\epsilon \to 0$. At tolerance 0.1, we found 227 near-orthogonal pairs per 30-ray configuration on average. However, as $\epsilon$ was tightened toward exact orthogonality, all near-orthogonal pairs evaporated: zero exact orthogonal pairs survived.

Total runtime across all four methods: approximately 12 minutes.

\textbf{Key insight: exact orthogonality is measure-zero.} In the space of unit vectors in $\R^3$ or $\C^3$, the condition $u \cdot v = 0$ defines a measure-zero set. There is no gradient signal connecting near-orthogonal configurations to exact ones: the ``number of exact orthogonal pairs'' objective is constant (zero) almost everywhere in continuous configuration space, jumping discontinuously only on a set of measure zero. Continuous optimization is \emph{categorically} unable to find KS sets. Any successful search must exploit algebraic structure---number-ring coordinates, symmetry groups, or other algebraic constraints that force exact orthogonality by construction.

\textbf{Interpretation.} The failure of all four numerical methods is not independent evidence against the existence of a sub-31 KS set. A sub-31 set could exist with coordinates in an untested algebraic field. The negative control confirms that if such a set exists, it must have algebraic structure, and further search should focus on identifying algebraic fields with untested cancellation mechanisms.

%%----------------------------------------------------------------------
\section{Criticality and deletion-minimality}
\label{sec:criticality}
%%----------------------------------------------------------------------

A KS set $S$ is \emph{deletion-minimal} if removing any single ray from $S$ produces a colorable set. Deletion-minimality implies that $S$ has no proper KS subset, since KS-uncolorability is hereditary under deletion.

\textbf{CK-31 is deletion-minimal.} We tested all 31 single-ray removals from CK-31 exhaustively via SAT. In each case, the resulting 30-ray set is colorable. This confirms that CK-31 contains no proper KS subset: \emph{every} proper subset of CK-31 is colorable. This result was verified independently for all six distinct minimal 31-sets found in the integer pool; each is deletion-minimal.

\textbf{Extended criticality testing.} As a computational diagnostic, we performed extended criticality testing by removing $k$ rays for $k = 1, \ldots, 12$ and testing colorability of the remaining $31 - k$ rays. Results are summarized in Table~\ref{tab:criticality}.

\begin{table}[t]
\centering
\caption{Extended criticality of CK-31. For $k \leq 8$, all $\binom{31}{k}$ subsets tested exhaustively; for $k = 9$--$12$, $10^6$ random subsets sampled per $k$. All subsets colorable.}
\label{tab:criticality}
\begin{tabular}{rrrl}
\toprule
$k$ removed & Remain & Subsets tested & Method \\
\midrule
1--4 & 30--27 & 36{,}456 & Exhaustive \\
5 & 26 & 169{,}911 & Exhaustive \\
6 & 25 & 736{,}281 & Exhaustive \\
7 & 24 & 2{,}629{,}575 & Exhaustive \\
8 & 23 & 7{,}888{,}725 & Exhaustive \\
9--12 & 22--19 & $4 \times 10^6$ & Sampled \\
\bottomrule
\end{tabular}
\end{table}

Total: 11{,}460{,}647 exhaustive SAT checks plus $4 \times 10^6$ sampled checks. \emph{All subsets are colorable.}

\textbf{Significance at the lower-bound threshold.} At $k = 7$, the remaining 24 rays equal the Li--Bright--Ganesh lower bound~\cite{LiBrightGanesh2024}. All $\binom{31}{7} = 2{,}629{,}575$ subsets of CK-31 of size 24 are colorable, providing explicit confirmation at the lower-bound threshold. (This already follows from deletion-minimality, but the exhaustive check provides a concrete certificate at exactly the boundary.)

\textbf{Implication.} The extended criticality data rules out all sub-configurations of CK-31 as potential smaller KS sets. Any sub-31 KS set must use rays \emph{not} in CK-31---drawn from a different part of the integer pool, from a different algebraic pool entirely, or from an untested number ring. Together with the vertex-merging unrealizability result (Section~\ref{sec:merging}), this closes off two natural routes to a sub-31 KS set within the integer pool: sub-configurations of CK-31 and combinatorial modifications of CK-31's orthogonality graph.

%%----------------------------------------------------------------------
\section{Triad density analysis}
\label{sec:density}
%%----------------------------------------------------------------------

The previous sections establish negative results for specific constructions. A complementary question is: how rare are KS sets within their own pools? If KS sets are extremely rare at sizes just above the pool minimum, sub-threshold existence is unlikely, though density arguments cannot constitute a proof of optimality.

\textbf{Method.} For each pool, we sampled 500 uniform random subsets at each size $n$, tested each for KS-uncolorability via SAT, and recorded the fraction that are uncolorable. We use the rule-of-three to set upper bounds on prevalence when no uncolorable sets are observed: if $k$ samples are drawn with no successes, the 95\% upper bound is $p \lesssim 3/k$.

\textbf{Integer pool (49 rays).} No KS-uncolorable subset is observed below $n = 42$ out of 49 rays ($p \lesssim 0.006$ at 95\% confidence). The gap between the SAT minimum of~31 and the density threshold of~42 reflects the extreme rarity of KS sets even within pools known to contain them.

\textbf{Eisenstein pool (57 rays).} No KS-uncolorable subset below $n = 48$ out of 57. The Eisenstein pool has a \emph{higher} SAT minimum (33 vs.\ 31) despite $\C^3$ providing additional orthogonality freedom, demonstrating that working in $\C^3$ does not automatically improve the minimum size.

\textbf{Heegner-7 pool (145 rays).} No KS-uncolorable subset at any sampled size up to $n = 75$ out of 145. This pool has the largest gap between pool size and SAT minimum (43 vs.\ 145), consistent with its sparser orthogonality structure.

\textbf{Pair density of CK-31.} Among all tested constructions, CK-31 achieves the highest pair density: 15.3\% of all $\binom{31}{2} = 465$ pairs are orthogonal.

\textbf{Key insight: density alone does not determine uncolorability.} KS-uncolorable sets are extraordinarily rare even within their own pools. The Eisenstein pool (complex coordinates) has a higher density threshold than the integer pool despite having additional orthogonality freedom. The boundary between colorable and uncolorable is determined by the \emph{topological interlocking} of the constraint hypergraph, not by any simple numerical density measure.

Random rays in $\R^3$ or $\C^3$, drawn uniformly from the unit sphere without reference to any algebraic pool, produce zero orthogonal pairs at any tested size $n \leq 50$ across 500 samples, confirming from a different angle that algebraic structure is essential.

\textbf{Implications for proof strategies.} This analysis rules out density-based approaches to proving optimality. A proof that no sub-31 KS set exists cannot rely on counting orthogonal pairs or triads. The precise topological structure of the constraint hypergraph must be the object of study.

%%----------------------------------------------------------------------
\section{Graph perturbation search}
\label{sec:perturbation}
%%----------------------------------------------------------------------

\subsection{Motivation}

Strategies~1--6 all search within or near established algebraic pools. If a sub-31 KS set exists with an orthogonality graph not arising from any known algebraic construction---one that is in some sense ``irregular''---those strategies would miss it entirely. Strategy~4 (numerical optimization) attempts to escape algebraic pools, but exact orthogonality is measure-zero in continuous space, so random ray replacement quickly loses all triad structure. A more targeted approach is needed: one that searches the space of abstract orthogonality hypergraphs directly, without requiring coordinates, then tests realizability only for uncolorable candidates.

\subsection{Method}

We start from CK-31's orthogonality graph: 31 vertices, 71 edges, 17 triads. Each trial applies between 1 and 5 perturbations drawn from five operations:
\begin{enumerate}
\item \emph{Remove vertex} (weight~3): delete one vertex and all incident edges; remap remaining vertices.
\item \emph{Swap triad} (weight~3): remove one triad and add a randomly chosen new 3-clique.
\item \emph{Merge vertices} (weight~2): identify two non-adjacent vertices into one, combining their neighborhoods.
\item \emph{Add pair} (weight~1): add an orthogonality edge between two non-adjacent vertices.
\item \emph{Remove pair} (weight~1): remove an orthogonality edge not belonging to any triad.
\end{enumerate}

After perturbation, we test KS-uncolorability via SAT (Glucose4). If UNSAT and $n < 31$, we test $\R^3$ realizability: minimize $\sum_{(i,j)\in E}(v_i\cdot v_j)^2$ for unit vectors $v_i \in \R^3$ using L-BFGS-B with analytical gradients and 10--15 random initializations, tolerance $10^{-10}$.

\subsection{Results}

We ran 500,000 trials over 12.73 hours. Table~\ref{tab:perturbation} summarizes the outcomes.

\begin{table}[t]
\centering
\caption{Graph perturbation search results (500,000 trials, 12.73 hours).}
\label{tab:perturbation}
\begin{tabular}{lr}
\toprule
Metric & Count \\
\midrule
Valid perturbations & 500,000 \\
Still KS-uncolorable & 79,850 (16.0\%) \\
Sub-31 and uncolorable & 61,702 (12.3\%) \\
Near-realizable (residual $< 0.1$) & 1,904 \\
Realizable (residual $< 10^{-10}$) & 0 \\
Best residual & $3.8\times 10^{-2}$ \\
Best configuration & $n=30$, $t=17$, $p=69$ \\
\bottomrule
\end{tabular}
\end{table}

\subsection{Interpretation}

\textbf{Abstract sub-31 KS-uncolorable hypergraphs are abundant.} Roughly 12.3\% of all valid perturbations produce abstract hypergraphs that have fewer than 31 vertices and are KS-uncolorable. The combinatorial requirement for a sub-31 KS set is not the bottleneck.

\textbf{The geometric requirement is the bottleneck.} Of the 61,702 abstract sub-31 candidates, numerical optimization (L-BFGS-B, 10--15 random starts per candidate) finds no $\R^3$ realization for any of them. The 1,904 near-realizable cases (residual $< 0.1$) represent the closest the optimizer reaches, and the best residual across all candidates is $3.8\times 10^{-2}$---more than six orders of magnitude above the acceptance threshold of $10^{-10}$. While numerical optimization cannot prove unrealizability, the systematic failure across 61,702 independent instances with consistently large residuals constitutes strong evidence of a structural obstruction rather than an optimization artifact.

\textbf{The best residual signals a structural gap, not a near-miss.} A residual of $3.8\times 10^{-2}$ reflects an incompatibility between the abstract graph's structure and three-dimensional real geometry. This is qualitatively different from a near-miss where the residual is small but optimization is slow.

\textbf{The best configuration closely mirrors CK-31.} The configuration achieving the best residual has $n = 30$, $t = 17$, $p = 69$. CK-31 has $n = 31$, $t = 17$, $p = 71$. Same number of triads, two fewer pairs, one fewer vertex---but the geometry does not give way.

\begin{remark}
The 61,702 abstract sub-31 KS-uncolorable hypergraphs are combinatorial objects: they exist as hypergraphs but no $\R^3$ realization has been found for any of them. Their existence shows that abstract KS-uncolorability is achievable below 31 vertices; the obstruction to a sub-31 KS \emph{set} (i.e., a realized configuration of rays) appears to lie in the interaction of uncolorability requirements with three-dimensional Euclidean geometry.
\end{remark}

%%----------------------------------------------------------------------
\section{Vertex merging and the realizability barrier}
\label{sec:merging}
%%----------------------------------------------------------------------

\subsection{The vertex merging construction}

We generate 30-vertex candidate KS graphs by \emph{vertex merging}: identify two non-adjacent vertices $u$ and $v$ into a single vertex $w$ whose neighborhood is $N(w) = N(u) \cup N(v) \setminus \{u, v\}$.

\begin{proposition}[Merge preserves KS-uncolorability]
\label{prop:merge}
If $G$ is the orthogonality graph of a KS-uncolorable ray set, and $G'$ is obtained from $G$ by merging two non-adjacent vertices $u$ and $v$, then $G'$ is KS-uncolorable.
\end{proposition}

\begin{proof}
Suppose $G'$ admits a KS coloring $f'$. Define $f: V(G) \to \{0,1\}$ by $f(u) = f(v) = f'(w)$ and $f(x) = f'(x)$ for all other vertices. Since $u$ and $v$ are non-adjacent in $G$, pairwise exclusion is not imposed between them. Every triad and pair constraint in $G$ involving at most one of $\{u,v\}$ is satisfied by $f'$. Constraints involving both $u$ and $v$ cannot arise from triads (since $u \not\perp v$). Thus $f$ is a valid KS coloring of $G$, contradicting its uncolorability.
\end{proof}

\subsection{Applying merging to CK-31}

CK-31 has $\binom{31}{2} - 71 = 394$ non-adjacent pairs. By Proposition~\ref{prop:merge}, all 394 merges produce abstract 30-vertex KS-uncolorable graphs with 17--21 triads and 70--71 pairs.

\subsection{Realizability testing: three approaches}

\textbf{Z3/SMT.} We encoded $\R^3$ realizability as polynomial constraints and submitted each instance to Z3~\cite{Z3}. Result: ``unknown'' on all 394. Intractable at this scale.

\textbf{Finite pool SAT.} Assign each of the 30 abstract vertices to one of the 49 integer pool rays, subject to injectivity, orthogonality preservation, and non-adjacency constraints. All 394 instances UNSAT in under 1\,s each. Definitively rules out realizability within the $\{0,\pm1,\pm2\}$ pool.

\textbf{Numerical optimization.} L-BFGS-B with 10--15 random starts per instance. No realization found; best residuals range from 0.01 to 0.15.

\subsection{The realizability barrier}

\begin{observation}[Realizability barrier]
\label{obs:barrier}
Two independent lines of evidence converge:
\begin{enumerate}[label=(\roman*)]
\item \emph{Vertex merging}: all 394 combinatorially valid 30-vertex KS graphs derived from CK-31 are proved unrealizable within the $\{0,\pm1,\pm2\}$ integer pool (finite pool SAT, all UNSAT in $< 1$\,s each). This is an exact result within the pool; realizability in other pools or in $\R^3$ generally remains open.
\item \emph{Graph perturbation} (Section~\ref{sec:perturbation}): among 61,702 abstract sub-31 KS-uncolorable hypergraphs, numerical optimization finds no $\R^3$ realization (best residual $3.8 \times 10^{-2}$). This is strong numerical evidence but not a proof of unrealizability.
\end{enumerate}
In total, at least 62,096 distinct abstract sub-31 KS-uncolorable hypergraphs have been tested; no realization has been found by any method.
\end{observation}

The evidence points to a \emph{realizability gap}: abstract KS-uncolorable hypergraphs with $n < 31$ exist in abundance, but embedding them as rays in $\R^3$ appears to be the core obstruction. The 394 merged graphs are definitively excluded from the integer pool; the broader claim that no sub-31 abstract KS hypergraph is $\R^3$-realizable is supported by extensive numerical evidence but remains unproved.

\begin{remark}
The finite pool SAT encoding resolves realizability instances in under 1\,s that defeat Z3. It is directly applicable as a ``quick accept'' screen in the Li--Bright--Ganesh pipeline~\cite{LiBrightGanesh2024}, which is bottlenecked by Z3 above order~20.
\end{remark}

\subsection{Open questions}

The 394 merged graphs were tested only within the integer pool. Open: (1)~realizability in other algebraic pools; (2)~realizability in~$\C^3$; (3)~isomorphism classes among the 394 graphs; (4)~vertex merging applied to other pools (e.g., 33-vector Eisenstein set).

%%----------------------------------------------------------------------
\section{Toward a proof of optimality}
\label{sec:proof}
%%----------------------------------------------------------------------

The computational evidence is strong: no sub-31 KS set has been found by any of seven independent methods, and the realizability barrier shows that abstract sub-31 KS hypergraphs are combinatorially plentiful but geometrically excluded. We frame the following as structural insights and candidate proof directions, not as partial proofs.

\subsection{Constraint budget}\label{sec:budget}

Suppose $S$ is a minimal KS set in~$\R^3$ with $|S| = n$. Let $t = |\mathcal{T}|$ denote the number of triads and $p = |E|$ the number of orthogonal pairs. Define $C = t + p$.

For CK-31: $n = 31$, $t = 17$, $p = 71$, $C = 88$, $C/n = 2.839$. The 71 pairs decompose into $51$ triad-implied pairs ($17 \times 3$) and $20$ extra pairs not in any triad.

\begin{observation}[Extra pairs are essential]
\label{obs:extra}
Every one of CK-31's 20 extra pairs is essential: removing any single extra pair makes the formula satisfiable. The 51 triad-implied pairs alone are satisfiable. CK-31 requires all 20 extra pairs simultaneously.
\end{observation}

CK-31 is \emph{maximally tight}: its uncolorability is spread across the entire constraint structure with no slack.

\subsection{Why counting arguments fail}\label{sec:counting}

\textbf{Degree-1 vectors are essential.} A naive intuition suggests every vector should participate in $\geq 2$ triads. This is false: 18 of 31 CK-31 vectors participate in only one triad. These degree-1 vectors are essential via extra pair constraints, not triads.

\begin{center}
\begin{tabular}{cc}
\toprule
Triads containing $v$ & Count \\
\midrule
1 & 18 \\
2 & 8 \\
3 & 3 \\
4 & 2 \\
\bottomrule
\end{tabular}
\end{center}

\textbf{High density does not imply uncolorability.} The densest 30-vertex subsets of the integer pool achieve $C/n$ up to 5.33---far above CK-31's 2.839---yet all are colorable. The precise topological interlocking, not density, determines uncolorability.

\subsection{The plane matching property}\label{sec:planematching}

\begin{definition}
For a ray $v \in \R^3$, the perpendicular plane $v^\perp$ is 2-dimensional. A triad $\{v, u, w\}$ places $u$ and $w$ in $v^\perp$ as an orthogonal pair.
\end{definition}

\begin{observation}[Plane matching property]
\label{obs:planematching}
In $\R^3$, if $v$ is contained in $k$ triads, the $2k$ triad neighbors of $v$ form $k$ orthogonal pairs (a matching) within $v^\perp$. This is verified for all 31 vectors of CK-31.
\end{observation}

The non-triviality of this constraint lies in its consequences for abstract hypergraph realizability. In $\R^3$, the 2-dimensionality of $v^\perp$ means that a vertex of triad-degree $k$ forces its $2k$ triad neighbors into $k$ orthogonal pairs within a single plane---a matching constraint that severely limits the number of triads that can share a vertex (at most $\lfloor (n-1)/2 \rfloor$ triads through any single ray, and each imposes a rigid directional relationship on its partners). In dimensions $d \geq 4$, $v^\perp$ has dimension $\geq 3$, orthogonal pairs within it need not form a matching, and the constraint vanishes. This dimensional restriction is a candidate explanation for why the minimum KS size in $\R^3$ is substantially larger than in $\R^4$ (where 18--20 vectors suffice).

The plane matching property is not itself a proof of optimality---it is an $\R^3$-specific structural constraint that any realized KS set must satisfy. The open question is whether it, combined with other geometric constraints, suffices to exclude all abstract KS-uncolorable hypergraphs below 31 vertices.

\begin{conjecture}[Local-global realizability gap]
\label{conj:planematching}
There exist abstract KS-uncolorable hypergraphs on $n < 31$ vertices whose local triad structure at each vertex is individually compatible with the plane matching property in~$\R^3$, but for which no globally consistent assignment of rays exists. That is, the obstruction to realizability is global, not detectable by local checks alone.
\end{conjecture}

\subsection{What a proof would require}\label{sec:proofstrategy}

A proof of optimality cannot rely on counting arguments. It must connect $\R^3$ geometry (which constrains realizable orthogonality graphs via the plane matching property) with the combinatorial structure required for uncolorability.

The Li--Bright--Ganesh approach~\cite{LiBrightGanesh2024} enumerates abstract hypergraphs bottom-up; our approach identifies geometric constraints top-down. A proof likely needs both: LBG-style enumeration (possibly accelerated by our finite SAT encoding) plus geometric constraints (plane matching, rigidity).

\begin{conjecture}[Optimality of CK-31]
\label{conj:optimality}
No KS set in $\R^3$ has fewer than 31 vectors. Equivalently, every abstract KS-uncolorable hypergraph on at most 30 vertices is unrealizable in~$\R^3$.
\end{conjecture}

Our computational evidence is consistent with this conjecture: 62,096 abstract sub-31 KS hypergraphs tested, zero realized.

%%----------------------------------------------------------------------
\section{Discussion: the landscape of approaches}
\label{sec:discussion}
%%----------------------------------------------------------------------

\subsection{Taxonomy of approaches}

The search for sub-31 KS sets has been pursued through four broad families of
methods, each with characteristic strengths and limitations.

\textbf{Algebraic search} (Strategies 1--3 and 6; pool construction, alphabet
expansion, and cross-pool mixing). The algebraic approach builds ray pools from
coordinate alphabets drawn from specific number rings, then applies SAT-based
uncolorability testing and OCUS extraction to find or rule out KS sets within
each pool. Its primary strength is exhaustiveness within a tested field: every
subset of a finite pool is implicitly checked, and OCUS certification provides
a formally verifiable guarantee that no smaller KS set exists in that pool.
Results are SAT-certifiable and reproducible. The principal limitation is scope:
there are infinitely many number fields, and untested fields may harbor novel
cancellation identities that enable new KS constructions.
Our contribution across this approach: six pools from distinct number rings are
fully characterized; the modulus-2 boundary is identified as the algebraic
obstruction separating 29-alphabet expansions (all failing) from 31-achieving
constructions; 30 expanded alphabets are tested without yielding sub-31 sets;
and OCUS provides a certified lower bound of 31 for the integer pool.

\textbf{Combinatorial enumeration}
\cite{LiBrightGanesh2024,KirchwegarPeitlSzeider2023}
(abstract hypergraph generation followed by realizability checking).
The combinatorial approach generates abstract KS-uncolorable hypergraphs
directly, without reference to a specific coordinate field, and then asks
whether each abstract graph can be realized by unit vectors in $\R^3$.
Its theoretical strength is completeness: any KS set that exists will eventually
appear as a node in the enumeration. The practical limitation is the realizability
bottleneck. Above approximately order~20, the Z3 SMT solver returns ``unknown''
on the nonlinear real arithmetic queries that realizability requires, so the
pipeline stalls even when abstract candidates are abundant.
Our finite pool SAT encoding (Section~\ref{sec:methods}) resolves instances
in under 1\,s that defeat Z3, and is directly applicable as a ``quick accept''
screen in the LBG pipeline: a combinatorial candidate can be tested for
membership in any of our six algebraic pools before an expensive Z3 query is
invoked. The 61,702 abstract sub-31 candidates found in our perturbation search
(Section~\ref{sec:perturbation}) provide a concrete input set for such screening.

\textbf{Numerical and heuristic search} (Strategies 4 and 7; simulated
annealing, perturbation, and graph perturbation). Heuristic methods explore
large search spaces without algebraic constraints, probing for KS sets that
might lie outside any specific algebraic pool. Their strength is speed and
breadth: 500,000 perturbation trials can be completed where an exhaustive
algebraic search would be infeasible. The fundamental limitation is the absence
of completeness guarantees: heuristic methods can find KS sets but cannot rule
them out. Our contribution is the quantification of the realizability barrier.
The perturbation search generated 61,702 distinct abstract KS-uncolorable
graphs of order 28--30, yet none was realizable in $\R^3$ under any of three
complementary realizability tests. This systematic failure is not a null result:
it establishes that the obstruction to sub-31 KS sets is not merely a lack of
uncolorable graphs (which exist in abundance) but specifically the failure of
those graphs to be embeddable in three-dimensional Euclidean space.

\textbf{Structural and proof-theoretic methods} (Section~\ref{sec:proof};
Trandafir--Cabello rigidity \cite{TrandafirCabello2025}; constraint budget
analysis; plane matching). Structural methods seek not merely to test whether
sub-31 KS sets exist but to prove or disprove their existence via geometric
or combinatorial arguments. This is the only approach capable of yielding a
definitive answer to the central question. The limitation is that no such proof
currently exists. Our contributions include: identifying why naive counting
arguments fail to close the gap; formulating the plane matching property as a
candidate geometric constraint that any sub-31 set would have to violate; and
the constraint budget analysis showing that $\CK$ is maximally tight in the
sense that its orthogonality constraints leave no slack for simplification.

\subsection{Connections to published work}

\textbf{Trandafir--Cabello rigidity} \cite{TrandafirCabello2025}.
Trandafir and Cabello proved that $\CK$ is unique up to unitary equivalence:
any KS set with exactly the same orthogonality graph as $\CK$ is unitarily
equivalent to it. Any sub-31 KS set must therefore have a different graph.
Combined with our result in \cite{Kernaghan2026algebraic} that all tested
algebraic constructions achieving 31 rays produce the same graph as $\CK$, this
implies that no tested construction yields an alternative 31-vertex graph that
might be perturbed downward. Rigidity and algebraic uniformity together suggest
that any sub-31 route would require a genuinely new construction principle.

\textbf{Cortez--Morales-Reyes $N(S)$ invariant} \cite{CortezMoralesReyes2022}.
Cortez and Morales-Reyes introduced the invariant $N(S) = \mathrm{lcm}\{\|v\|^2
: v \in S\}$ for a KS set $S$ with integer coordinates, and proved that $S$
must be defined over the ring $\Z[1/N(S)]$. For $\CK$, $N = 30$, corresponding
to the ring $\Z[1/30]$. Our Eisenstein pool has $N = 6$, precisely the minimal
ring $\Z[1/6]$ predicted by their theorem for Eisenstein-coordinate
constructions. This provides an independent algebraic classification: different
KS islands inhabit distinct arithmetic strata, and the $N(S)$ invariant
constrains which rings need to be searched for each target size.

\textbf{SI-C closure convergence.}
Trandafir and Cabello's 97-ray SI-C closure of $\CK$ \cite{TrandafirCabello2025}
and our 49-ray integer pool overlap in 37 rays. After a second round of closure
(producing 1741 rays), our entire integer pool is contained in the closure.
This convergence of the top-down (closure from $\CK$) and bottom-up (alphabet
enumeration) approaches suggests that the algebraically accessible landscape
of integer-coordinate KS sets is complete, or very nearly so. Rays outside this
landscape would require coordinates from genuinely different number fields.

\subsection{Assessment}

The two main open problems---improving the lower bound from 24 and proving that
31 is optimal---have very different tractability profiles. Improving the lower
bound to 28 or 29 may be achievable with current technology: our finite pool
SAT encoding accelerates the LBG pipeline's realizability bottleneck, and the
modulus-2 boundary and plane matching property provide additional constraints
that could prune the combinatorial enumeration. Even a modest improvement to 28
or 29 would be significant and publishable. Proving that 31 is optimal is harder.
A proof would need to combine geometric realizability obstructions with
combinatorial uncolorability in a way that no current technique accomplishes.

Taken together, the results of this series define where a hypothetical sub-31
KS set could hide: it would need to arise from a number field not yet tested,
exhibit a novel cancellation identity that our alphabet expansion strategy missed,
produce an abstract graph that is both uncolorable and geometrically realizable
in $\R^3$, and have no realization in any of the six pools studied here.
Each condition is individually possible; their simultaneous satisfaction becomes
increasingly implausible in light of the evidence assembled.

%%----------------------------------------------------------------------
\section{Conclusion and open problems}
\label{sec:conclusion}
%%----------------------------------------------------------------------

We have described seven complementary strategies---algebraic pool construction
and OCUS certification, alphabet expansion, cross-pool mixing, simulated
annealing, modulus-based ring analysis, graph perturbation, and structural
proof-theoretic analysis---applied systematically to six algebraic ray pools
drawn from distinct number rings. In no case was a KS set with fewer than 31
rays found. The central finding is the identification of the \emph{realizability
barrier} as the primary obstruction: abstract KS-uncolorable graphs of order
28--30 exist in abundance (61,702 distinct candidates from the perturbation
search alone), but $\R^3$ geometry prevents their realization. The difficulty
of closing the 24--31 gap is not a shortage of combinatorial candidates but
the failure of those candidates to be embeddable in three-dimensional Euclidean
space with the required orthogonality structure.

The following problems remain open.

\begin{enumerate}

\item \textbf{Is 31 optimal?} The central open problem of the field. A proof
would require combining geometric realizability constraints with combinatorial
uncolorability in a single argument---a goal that the constraint budget
analysis and plane matching results of Section~\ref{sec:proof} approach but do
not reach. \textit{(Hard; no current technique is sufficient.)}

\item \textbf{Can the lower bound be improved beyond 24?} The finite pool SAT
encoding developed here resolves realizability queries in under 1\,s that defeat
Z3, and is immediately applicable as a screening layer in the LBG combinatorial
enumeration pipeline \cite{LiBrightGanesh2024}. Geometric constraints from the
plane matching property and Trandafir--Cabello rigidity \cite{TrandafirCabello2025}
may further prune the search space. An improvement to 28 or 29 may be achievable
with current SAT technology combined with these algebraic and geometric
constraints. \textit{(Tractable; the most promising near-term target.)}

\item \textbf{Does the modulus-2 boundary hold for all number fields?} We have
verified that every tested alphabet lacking the coordinate $\pm 2$ fails to
produce a KS-uncolorable pool, while every alphabet containing $\pm 2$ succeeds.
This has been established for quadratic, cyclotomic, and Pisot algebraic
families. Whether the boundary holds for higher-degree algebraic extensions or
transcendental coordinates is untested. \textit{(Open; requires new
computational experiments or a ring-theoretic argument.)}

\item \textbf{Are any of the 394 merged 30-vertex graphs realizable in
$\mathbb{C}^3$?} Section~\ref{sec:merging} identifies 394 uncolorable
30-vertex graphs produced by vertex merging, none realizable in $\R^3$ under
our tests. Testing realizability over $\mathbb{C}^3$ is a straightforward
extension of the numerical optimization method (replacing real dot products
with Hermitian inner products) not yet carried out. A positive result would
be dramatic: a 30-ray KS set over the complex numbers in dimension three.
\textit{(Computable; an accessible near-term experiment.)}

\item \textbf{Does the plane matching property yield combinatorial constraints
at $n = 30$?} The plane matching observation establishes a structural
regularity of all known KS sets at the CK-31 scale. A combinatorial proof
that no uncolorable 3-uniform hypergraph with the matching property exists
at order 30 would constitute a direct proof of optimality. \textit{(Research
question; requires a new combinatorial argument.)}

\item \textbf{Can merge saturation be proved for all minimal KS sets?} Vertex
merging is verified to preserve KS-uncolorability across all 3,756 tested
merges over the six known KS islands, with 100\% preservation. This empirical
saturation---that every merge of a minimal KS set yields another KS
set---suggests a structural characterization of minimal KS sets that, if
proved, would substantially constrain the space of candidates. \textit{(Open
conjecture; well-supported computationally.)}

\end{enumerate}

The 24--31 gap remains one of the central open problems in Kochen--Specker
theory. The results of this series have mapped the landscape of approaches
more completely than previously available and have identified the realizability
barrier as the specific geometric obstruction that the search encounters.
The evidence assembled here---the algebraic uniformity of all tested
constructions, the failure of 61,702 abstract candidates to realize in $\R^3$,
the maximally tight constraint budget of $\CK$, and the convergence of
top-down and bottom-up algebraic explorations---points consistently toward
the conjecture that 31 is optimal.
A proof remains elusive.

%%----------------------------------------------------------------------
\begin{acknowledgments}
The author thanks Ad\'an Cabello for valuable correspondence on KS set rigidity
and for offering to sponsor the arXiv submission.
The author acknowledges the use of Claude (Anthropic) for computational assistance.
SAT solving used PySAT with the Glucose4 solver.
Code, ray/triad lists, and SAT instances are available at
\url{https://github.com/michaelkernaghan/contextuality}.
\end{acknowledgments}

%%----------------------------------------------------------------------
\begin{thebibliography}{99}

\bibitem{KochenSpecker1967}
S.~Kochen and E.~P.~Specker, ``The problem of hidden variables in quantum
mechanics,'' \textit{J. Math. Mech.} \textbf{17}, 59--87 (1967).

\bibitem{Peres1991}
A.~Peres, ``Two simple proofs of the Kochen--Specker theorem,''
\textit{J. Phys. A: Math. Gen.} \textbf{24}, L175--L178 (1991).

\bibitem{Peres1993}
A.~Peres, \textit{Quantum Theory: Concepts and Methods}
(Kluwer Academic Publishers, 1993).

\bibitem{ConwayKochen}
J.~H.~Conway and S.~Kochen, reported in A.~Peres,
\textit{Quantum Theory: Concepts and Methods} (Kluwer, 1993).

\bibitem{Kernaghan1994}
M.~Kernaghan, ``Bell--Kochen--Specker theorem for 20 vectors,''
\textit{J. Phys. A: Math. Gen.} \textbf{27}, L829--L830 (1994).

\bibitem{KernaghanPeres1995}
M.~Kernaghan and A.~Peres, ``Kochen--Specker theorem for eight-dimensional
space,'' \textit{Phys. Lett. A} \textbf{198}, 1--5 (1995).

\bibitem{CabelloEstebaranzGarcia1996}
A.~Cabello, J.~M.~Estebaranz, and G.~Garc\'{\i}a-Alcaine,
``Bell--Kochen--Specker theorem: A proof with 18 vectors,''
\textit{Phys. Lett. A} \textbf{212}, 183--187 (1996).

\bibitem{Pavicic2004}
M.~Pavicic, J.-P.~Merlet, B.~McKay, and N.~D.~Megill,
``Kochen--Specker vectors,''
\textit{J. Phys. A: Math. Gen.} \textbf{38}, 1577--1592 (2005).
arXiv:quant-ph/0409014.

\bibitem{ArendsOuaknineWampler2011}
F.~Arends, J.~Ouaknine, and C.~W.~Wampler,
``On searching for small Kochen--Specker vector systems,''
in \textit{Proceedings of the International Workshop on Graph-Theoretic
Concepts in Computer Science (WG 2011)}, LNCS 6986, 23--34 (2011).

\bibitem{Pavicic2019}
M.~Pavicic, N.~D.~Megill, B.~C.~Waegell, and P.~K.~Aravind,
``Automated generation of Kochen--Specker sets,''
\textit{Sci. Rep.} \textbf{9}, 6765 (2019). arXiv:1905.01009.

\bibitem{UijlenWesterbaan2016}
S.~Uijlen and B.~Westerbaan,
``A Kochen--Specker system has at least 22 vectors,''
\textit{New Generation Computing} \textbf{34}, 3--23 (2016).

\bibitem{AudemardSimon2018}
G.~Audemard and L.~Simon,
``On the Glucose SAT solver,''
\textit{Int. J. Artif. Intell. Tools} \textbf{27}, 1840001 (2018).

\bibitem{Z3}
L.~de~Moura and N.~Bj{\o}rner,
``Z3: An efficient SMT solver,''
in \textit{Proceedings of TACAS 2008}, LNCS 4963, 337--340 (2008).

\bibitem{CortezMoralesReyes2022}
V.~Cortez, R.~Morales, and E.~J.~Reyes,
``Minimal ring extensions of the integers exhibiting Kochen--Specker
contextuality,''
arXiv:2211.13216 [quant-ph] (2022, preprint).

\bibitem{Budroni2022}
C.~Budroni, A.~Cabello, O.~G\"uhne, M.~Kleinmann, and J.-\AA.~Larsson,
``Kochen--Specker contextuality,''
\textit{Rev.\ Mod.\ Phys.} \textbf{94}, 045007 (2022).

\bibitem{KirchwegarPeitlSzeider2023}
M.~Kirchweger, A.~Peitl, and S.~Szeider,
``Co-certificate learning with SAT modulo symmetries,''
in \textit{Proceedings of AAAI-23}, 4119--4126 (2023).

\bibitem{LiBrightGanesh2024}
J.~Li, C.~Bright, and V.~Ganesh,
``Leveraging SAT solvers for the minimum Kochen--Specker problem,''
in \textit{Proceedings of IJCAI-24}, 1944--1952 (2024).
arXiv:2306.13319.

\bibitem{Cabello2025simplest}
A.~Cabello,
``Simplest Kochen--Specker set,''
\textit{Phys.\ Rev.\ Lett.} (in press, 2025). arXiv:2508.07335.

\bibitem{TrandafirCabello2025}
S.~Trandafir and A.~Cabello,
``Two fundamental solutions to the rigid Kochen--Specker set problem and
the solution to the minimal Kochen--Specker set problem under one assumption,''
\textit{Phys.\ Rev.\ A} \textbf{111}, 052204 (2025). arXiv:2501.11640.

\bibitem{Kernaghan2026algebraic}
M.~Kernaghan,
``The algebraic landscape of Kochen--Specker sets in dimension three,''
arXiv:XXXX.XXXXX [quant-ph] (2026).

\end{thebibliography}

\end{document}
